\chapter{Protocolo de Medición Res. SRT 85/2012}
\label{sec:protocolo_srt}

En este capítulo se transcribe el cuadro oficial del \textbf{Protocolo de Medición de Ruido en el Ambiente Laboral}
conforme al formato exacto de la Resolución S.R.T. N° 85/2012, volcando los datos obtenidos en el ensayo con el minitorno Dremel.

\begin{table}[!ht]
  \centering
  \caption{Protocolo de Medición de Ruido en el Ambiente Laboral (Res. SRT 85/2012).}
  \label{tab.ej1:protocolo_srt_teorico}
  \resizebox{\textwidth}{!}{%
    \begin{tabular}{|c|c|c|c|c|c|c|c|c|c|c|}
      \hline
      \multicolumn{11}{|c|}{\textbf{PROTOCOLO DE MEDICIÓN DE RUIDO EN EL AMBIENTE LABORAL}} \\
      \hline
      \multicolumn{6}{|l|}{\textbf{Razón social:} Servicio Técnico Taller Electrodomésticos} & \multicolumn{5}{l|}{\textbf{C.U.I.T.:} 20-40191000-8} \\
      \hline
      \multicolumn{3}{|l|}{\textbf{Dirección:} Maestro M. López esq. Cruz Roja} & \multicolumn{3}{l|}{\textbf{Localidad:} Córdoba} & \multicolumn{2}{l|}{\textbf{C.P.:} 5000} & \multicolumn{3}{l|}{\textbf{Provincia:} Córdoba} \\
      \hline
      \multicolumn{11}{|c|}{\textbf{DATOS DE LA MEDICIÓN}} \\
      \hline
      & & & \textbf{Tiempo de} & \textbf{Tiempo de} & \textbf{Características} & \textbf{RUIDO DE IMPACTO} & \multicolumn{3}{c|}{\textbf{SONIDO CONTINUO o INTERMITENTE}} & \textbf{Cumple con} \\
      \cline{8-10}
      \textbf{Punto de} & \textbf{Sector} & \textbf{Puesto / Puesto} & \textbf{exposición} & \textbf{integración} & \textbf{generales del} & \textbf{Nivel pico $L_{C,pico}$} & \textbf{Nivel integrado} & \textbf{Suma de} & \textbf{Dosis} & \textbf{valores} \\
      \textbf{medición} & & \textbf{tipo / móvil} & \textbf{($T_e$, en hs)} & \textbf{(medición)} & \textbf{ruido a medir} & \textbf{(en dBC)} & \textbf{($L_{Aeq,Te}$ en dBA)} & \textbf{fracciones} & \textbf{(\%)} & \textbf{permitidos?} \\
      \hline
      1 & Taller Reparación & Operador Dremel & 8 & $306,7\text{ s}$ & Intermitente & 91,6 & 84,7 & 0,92 & 92,5 & \textbf{SÍ} \\
      \hline
    \end{tabular}%
  }
\end{table}

\section{Análisis de Datos y Mejoras a Realizar}

De acuerdo con la sección de conclusiones y recomendaciones del formulario oficial de la Res. SRT 85/2012, se vuelca el
siguiente diagnóstico técnico del puesto:

\begin{table}[!ht]
  \centering
  \caption{Conclusiones y Recomendaciones del Protocolo Res. SRT 85/2012.}
  \label{tab.ej1:protocolo_srt_conclusiones}
  \resizebox{\textwidth}{!}{%
    \begin{tabular}{|p{0.48\textwidth}|p{0.48\textwidth}|}
      \hline
      \multicolumn{2}{|c|}{\textbf{ANÁLISIS DE LOS DATOS Y MEJORAS A REALIZAR}} \\
      \hline
      \multicolumn{1}{|c|}{\textbf{Conclusiones}} & \multicolumn{1}{c|}{\textbf{Recomendaciones para adecuar el nivel de ruido}} \\
      \hline
      1. El puesto de \textbf{Operador de Minitorno Dremel} presenta un nivel sonoro continuo equivalente de $L_{Aeq} = 84,66\dBA$, cumpliendo con el límite máximo permisible de $85,00\dBA$ (Res. 295/03) para 8 horas. & 1. Proveer protectores auditivos de inserción o copa como medida preventiva recomendada debido al umbral elevado ($92,5\%$). \\
      2. La dosis diaria proyectada ($D\% = 92,5\%$) se encuentra dentro del margen tolerable ($<100\%$), pero operando en zona de precaución higiénica. & 2. Efectuar mantenimiento periódico en el minitorno y verificar la excentricidad de los accesorios de corte/pulido para reducir la vibración sonora. \\
      3. El tiempo máximo continuo de exposición sin superar la dosis legal es de $8\h \; 39\text{ min}$. & 3. Implementar controles audiométricos preventivos anuales para el personal del área de reparación. \\
      \hline
    \end{tabular}%
  }
\end{table}
